% ============================================================
%  §5  Discussion
% ============================================================
\label{sec:discussion}

\paragraph{When does evolution help?}
The results suggest a practical diagnostic rather than a universal taxonomy. If failures can be
fixed by changing the procedure the agent follows, skill evolution is the lower-cost first
choice. If the agent needs a different view of the image before the evidence is accessible,
tool evolution becomes necessary. A practitioner targeting a new benchmark can use a small pilot
run ($k{\approx}25$, $N{=}1$) to inspect which kinds of capabilities are generated and whether
they transfer beyond the triggering examples.

\paragraph{Limitations.}
Several limitations bound the scope of our conclusions.
\textit{(a) Short evolution horizon.}
We report results at $N{=}3$ iterations; the evolution dynamics suggest gains plateau
around $N{=}5$, but we have not evaluated whether longer runs (e.g.\ $N{>}10$) eventually
overfit the training distribution.
\textit{(b) Tool quality ceiling.}
Tool effectiveness is bounded by the VLM's Python coding ability.
For benchmarks requiring specialised computer vision operations beyond standard libraries
(e.g.\ learned feature extractors), the Generator may fail to produce valid tools, limiting
the perceptual gains.
\textit{(c) Benchmark scope.}
We evaluate on four main visual reasoning benchmarks plus GTA and VTool-R1-aligned splits.
Broader coverage across additional domains (e.g.\ video QA, 3D spatial reasoning, medical
imaging) is needed before making stronger claims about the generality of the proposed
cognitive/perceptual decision rule.
\textit{(d) API cost.}
Each training case incurs up to $\sim$15K tokens (attempt + failure analysis + generation +
validation), summing to $\sim$3M tokens per benchmark run.  While negligible compared to RL
training, this cost should be accounted for when comparing training-free methods.

\paragraph{Future directions.}
Three extensions merit investigation.
First, \emph{cross-benchmark skill transfer}: skills generated from ChartQA failures may
transfer to MathVista (both involve structured data reading); a shared skill library across
benchmarks could reduce per-benchmark cold-start cost.
Second, \emph{image generation as a reasoning tool}: rather than only processing existing images,
the agent could generate hypothetical visualisations (e.g.\ mentally redrawing a chart in a
clearer format) using image generation APIs, extending the ``think with images'' paradigm
\citep{vtoolr1_2025} to internally imagined percepts.
Third, \emph{online evolution}: rather than a fixed pre-deployment training phase, the agent
could evolve capabilities during deployment, adapting to the distribution shift of live queries.
